\section{Introduction}
\label{sec:introduction}

School mathematics is best at procedures: the National Research Council found students taught by traditional curricula more fluent in procedures than in understanding, strategy or reasoning~\cite{anonymous2001adding}. Solving a problem one has never seen is a different skill: understanding the problem, devising a plan, carrying it out and looking back~\cite{polya2014solve}. Olympiad-style problems call for exactly this skill. In our working definition, a problem is olympiad-style when its method is not immediately obvious to the solver, its solution rests on an idea rather than a procedure, it needs no knowledge beyond the child's grade, and, once the idea is found, the solution is short. Practising them needs a steady supply of problems, each new to the child, a little harder than the last, and still solvable.

Large language models can supply such problems without end, and a chat assistant puts them one message away. What a model writes, though, cannot go to a child unchecked. When models tutored college algebra, only \stat{literature}{gupta_dialogues_entirely_correct_percent}\pct{} of \stat{literature}{gupta_dialogues} tutoring dialogues were entirely correct~\cite{gupta2025beyond}, and teachers find some of the word problems models write unsolvable, wrongly answered or unfit for school, even from the strongest models~\cite{christ2024mathwell,christ2026edumath}. A wrong answer key turns a child's correct solution into a ``mistake'', and the child cannot tell which of the two is at fault.

We let different parties write and check. The chat's own model, which already talks to the child, writes each task. A small service that makes no model calls of its own proposes what the child should practise, hands the model what it needs, admits the task only after checking everything that can be checked mechanically, shows it without its answer and records the child's answer. We ask:
\begin{description}
\item[RQ1.] Which defects in a model-written task do the service's checks catch, and which do they miss?
\item[RQ2.] Does the learner model place a new child and keep practice in its target corridor, and how often does it declare a topic mastered when it is not?
\end{description}
We contribute an architecture in which the chat's model writes every task and a stateless service that calls no model checks it and seals its answer (Section~\ref{sec:system}); \statword{product}{check_codes} checks, centred on the model's own solver run twice under rotated option labels (Section~\ref{sec:verification}); a learner model for tasks used once, a hierarchical Elo model with a guessing floor on a designed ladder across grades \stat{product}{grade_min}--\stat{product}{grade_max}, with a trial series for placement, related to maximum likelihood (Section~\ref{sec:learner}); an offline evaluation with its negative results (Section~\ref{sec:evaluation}); and the code and \stat{product}{reference_tasks} reference tasks under the MIT license.
